\chapter{Histories Made of Events}
\label{ch:events}

\index{event}
\index{trace}
Every execution is represented by an event sequence. A deliberately small common
envelope makes histories comparable across protocols:

\begin{lstlisting}
structure Event (Actor Session Task Kind Payload : Type) where
  timeAt  : ClockTimestamp
  id      : String
  priorIds : List String
  session : Session
  task    : Task
  from    : Actor
  to      : Actor
  kind    : Kind
  payload : Payload

abbrev Trace (Actor Session Task Kind Payload : Type) :=
  List (Event Actor Session Task Kind Payload)
\end{lstlisting}

The fields have different jobs. \texttt{id} identifies a visible fact;
\texttt{priorIds} records zero or more immediate causal claims; \texttt{from} and
\texttt{to} name the parties; \texttt{(session, task)} correlates independently
progressing conversations; and \texttt{kind} supplies a stable discriminator.
The payload is protocol-specific and carries the parameters that flow from one
production to the next. \texttt{timeAt} is sampled from the model's shared monotonic
clock. Its differences measure elapsed time; it is not a civil-time date and its
origin is immaterial.

The causal links form a directed acyclic graph rather than a singly linked list.
An initial event has no prior identifiers. A fork is represented by several later
events that each name the same prior event. A join is represented by one event
whose \texttt{priorIds} names every branch on which it immediately depends.

\section{Compute everything from a fork/join event list}

Each event is one observation at one instant. Starting and completing are
different messages, so an event does not contain an interval. Durations are
derived by selecting the two boundary events and subtracting their timestamps.

\begin{lstlisting}
structure AnalysisEvent where
  id       : EventId
  priorIds : List EventId
  timeAt   : ClockTimestamp
  bytes    : ByteArray
  reveals  : List ValueRef

def elapsed (started completed : AnalysisEvent) : Duration :=
  completed.timeAt - started.timeAt
\end{lstlisting}

Consider the first order in \texttt{examples/bakery-events.jsonl}. Payment and
baking fork after the order is placed and join before the order is accepted:

\begin{center}
\small
\begin{tabular}{@{}llllrl@{}}
\toprule
Event & Immediate prior & \texttt{timeAt} & Meaning & Bytes & Reveals \\
\midrule
$e_0$ & $[]$ & $258726$ & order placed & 85 & order \\
$e_1$ & $[e_0]$ & $258921$ & payment requested & 71 & price \\
$e_2$ & $[e_0]$ & $259194$ & bake requested & 73 & product \\
$e_3$ & $[e_1]$ & $262219$ & payment authorized & 49 & authorization \\
$e_4$ & $[e_2]$ & $2055947$ & bake completed & 55 & completion \\
$e_5$ & $[e_3,e_4]$ & $2056288$ & order accepted & 57 & outcome \\
$e_6$ & $[e_5]$ & $2056291$ & delivery requested & 54 & delivery \\
$e_7$ & $[e_6]$ & $2283253$ & courier assigned & 40 & courier \\
$e_8$ & $[e_7]$ & $6946819$ & order delivered & 65 & outcome \\
\bottomrule
\end{tabular}
\end{center}

The list order is a serialization convenient for storage. Causal order is the
transitive closure of \texttt{priorIds}. Both $e_1$ and $e_2$ cite $e_0$, so they
fork. Event $e_5$ cites both $e_3$ and $e_4$, so it joins the branches. The
following quantities are folds or graph queries over this one list:

\begin{align*}
  \operatorname{paymentTime}(e_1,e_3) &= e_3.timeAt-e_1.timeAt = 3298,\\
  \operatorname{latency}(e_0,e_8) &= e_8.timeAt-e_0.timeAt = 6688093,\\
  \operatorname{wireBytes} &= 85+71+73+49+55+57+54+40+65 = 549.
\end{align*}

Two events are causally concurrent when neither is an ancestor of the other:

\begin{lstlisting}
def concurrent (events : List AnalysisEvent) (a b : EventId) : Bool :=
  !reachable events a b && !reachable events b a
\end{lstlisting}

The maximum concurrency permitted by the trace is the width of this partial
order: the size of its largest antichain. Here $\{e_1,e_2\}$ is an antichain, as
is $\{e_3,e_4\}$, and no three events are pairwise concurrent. Thus the maximal
causal concurrency is two. A timestamped point event has no active interval;
runtime concurrency requires separate start and completion messages. This is stronger than counting
adjacent list elements and different from measuring average runtime concurrency.

Timestamps and causal links also give elapsed time along a selected path:

\begin{lstlisting}
edgeTime(prior, event) := event.timeAt - prior.timeAt
elapsed(started, completed) := completed.timeAt - started.timeAt
\end{lstlisting}

The critical path is $e_0,e_2,e_4,e_5,e_6,e_7,e_8$, with elapsed time 6,688,093 ms.
The payment branch finishes at tick 262,219 and has 1,793,728 ms of slack before the join.
Queue occupancy, queued
bytes, active \texttt{(session,task)} keys, and memory pressure are computed by
adding enqueue/dequeue events to the same alphabet and folding them in clock
order.

The entire checked-in corpus contains \BakeryOrders{} orders and
\BakeryEvents{} events. Streaming that file produces the following model inputs;
the numbers in this table are generated LaTeX macros, not independently typed
constants:

\begin{center}
\begin{tabular}{@{}lr@{}}
\toprule
Reduction over \texttt{bakery-events.jsonl} & Observed value \\
\midrule
payment authorized / orders & \BakeryPaymentRatio \\
bake completed / orders & \BakeryBakeRatio \\
order accepted / orders & \BakeryAcceptRatio \\
delivered / orders & \BakeryDeliveredRatio \\
mean terminal latency & \BakeryMeanLatency{} ms \\
median terminal latency & \BakeryPFiftyLatency{} ms \\
95th-percentile terminal latency & \BakeryPNinetyFiveLatency{} ms \\
maximum simultaneously active orders & \BakeryMaxActive \\
\bottomrule
\end{tabular}
\end{center}

These ratios are suitable probability estimates for this corpus. Regeneration
with a different source stream changes the macros and therefore changes every
book example that consumes them.

Knowledge is another event reduction. For a value $v$, calculate the actors that
can derive it from initial knowledge and the bytes visible to them:

\begin{lstlisting}
def knowers (actors : List Actor) (v : ValueRef)
    (events : List AnalysisEvent) : List Actor :=
  actors.filter fun actor =>
    Derives (knowledgeFrom actor events) v
\end{lstlisting}

Encryption, signing, projection, key possession, and compromise contribute rules
to \texttt{Derives}. The primary result is the set \texttt{knowers v events}.
``Only these actors know $v$'' is a comparison between that computed set and an
allowed set, not a separate primitive called secrecy.

\section{Derived state}

A current session, an authenticated principal, or an outstanding request can often
be computed by folding the trace. Such a fact need not also exist as primitive
model state. This discipline has two benefits: evidence for a fact remains visible,
and the state representation does not multiply independently of the history that
justifies it.

Some implementations necessarily maintain private state. The model includes it
only when two executions with the same visible trace may make different future
behaviour legal. Even then, the first remedy is to expose the relevant distinction
as an event or a grammar parameter rather than copying an implementation object
graph into the specification.

\section{Time is more than a timestamp}

\index{time}
The central modelling problem is maintainable time evolution. We separate three
ideas that are often conflated:

\begin{itemize}
  \item sequence order: one visible event occurs before another;
  \item causality: an event extends or responds to another event;
  \item duration: an amount of physical or logical time elapses.
\end{itemize}

Grammar structure handles sequence, identifiers handle causal linkage, and typed
constraints handle duration. This keeps a change in timeout policy from rewriting
the protocol's entire behavioural shape.

For two events from the same \texttt{(session, task)}, elapsed time is defined when
the first event's \texttt{timeAt} value does not exceed the second's:

\begin{lstlisting}
def elapsed (start finish : Event) : Option Duration :=
  if start.session = finish.session &&
     start.task = finish.task &&
     start.timeAt <= finish.timeAt
  then some (finish.timeAt - start.timeAt)
  else none
\end{lstlisting}

Latency queries name their boundary event kinds and correlate them by
\texttt{(session, task)} before subtracting timestamps. Causality still comes
from grammar structure and \texttt{priorIds}: equal timestamps may occur, and a
timestamp alone does not establish that one event caused another.

\section{Round trip through bytes}

A generated scenario is not complete until it can cross the byte boundary in
both directions. \texttt{ScenarioEvent} in \texttt{Requirements.proto} carries
the event identity, complete \texttt{prior} list, correlation key, actors,
\texttt{timeAt}, and the selected protobuf atom. Encoding a \texttt{Scenario}
therefore produces bytes for both its causal envelope and payload. Decoding those
bytes must reconstruct the same ordered event list. The executable round-trip
check compares the decoded protobuf value with the generated value and separately
checks the reconstructed \texttt{prior} graph; loss of a fork or join fails the
build.
